\section{Python instruments and discontinuities: align the null, cutoff and returned set}
\label{sec:exec25-python-iv-rd}

A lending rule based on a credit-score cutoff and an encouragement that
changes take-up pose different causal questions. The first asks about a
local discontinuity under a continuity design; the second asks which
structural or instrument-induced contrast is identified. Their software
can share regression machinery without sharing a target or an uncertainty
claim. This comparison extends the retained linearmodels and pyfixest
kernel reviews with CausalPy IV/RD, rdrobust, rdd and ivmodels sources.
Every recommendation is conditional on the specified design and artifact.
\citep{exec25_pyrdiv_choice,exec22_pyiv_kernels_linearmodels,
exec22_pyiv_kernels_pyfixest}

\subsection{Preserve the IV target before improving computation}
The earlier linearmodels review distinguishes design and instrument rank
checks from rank of the projected design, instrument strength and exclusion.
It also distinguishes heteroskedastic covariance from a declared cluster or
calendar-dependence model. The retained pyfixest review follows weighted
two-stage regression, response-scale residuals and diagnostic distributions.
A first-stage relevance test or an overidentification test does not become
a coefficient-null confidence-set inversion because all are called IV
diagnostics. Those distinctions remain prerequisites for a bank adapter.
\citep{exec22_pyiv_kernels_linearmodels,exec22_pyiv_kernels_pyfixest,
exec22_pyiv_reference}

The newly inspected pyfixest delegate solves normal equations using a
selected NumPy or SciPy route. Least-squares and LSQR branches extract
the solution component, while the SciPy dense route declares a positive
matrix assumption. Thus a returned coefficient vector does not preserve
every available rank or iterative-termination diagnostic. The Rust
collinearity delegate uses an ordered Cholesky-style residual-diagonal
test on a Gram matrix. Column ordering and the absolute tolerance affect
which columns survive.~\citep{exec25_pyrdiv_pyfixest}

For an exact scaling example, a one-column Gram entry $10^{-8}$ lies below
a tolerance $10^{-6}$. Multiplying that column by 100 changes the entry to
$10^{-4}$, which exceeds the same tolerance although the algebraic rank
has not changed. Qualification must preserve units, ordering, removed
columns and tolerance. The inspected Rust error text mentions a positive
finite tolerance, but its local guard tests only whether the tolerance is
nonpositive. Nonfinite inputs therefore need an explicit boundary test and
complete caller review; this observation does not establish an observed
package failure.~\citep{exec25_pyrdiv_pyfixest,exec25_pyrdiv_choice}

The same source header says a reimplementation of fixest code has written
permission for MIT redistribution and names a permissions document. That
document is absent from the retained complete tree at this pin. The source
assertion remains recorded, but authentic permission bytes are not supplied
by repeating the assertion. File-level provenance is unresolved, without
a legal conclusion. The proposal keeps the numerical component as a
reference and defers new distribution or porting until the artifact-specific
record and native numerical contract are established.
\citep{exec25_pyrdiv_pyfixest,exec25_pyrdiv_choice}

\subsection{CausalPy's Bayesian IV route adds a model and prior contract}
The inherited CausalPy 1.0.0rc1 revision constructs separate outcome and
instrument formula matrices, computes deterministic OLS and two-stage
reference fits, and can use those fits to center default priors. Its IV
constructor does not sample; explicit build and fit steps construct the
model and sample the posterior. The inspected continuous-treatment backend
uses a joint multivariate-normal likelihood with an LKJ covariance
construction. Its binary-treatment branch uses a logistic latent variable,
a Bernoulli treatment likelihood and a conditional normal outcome model
with a clipped coupling parameter. These are additional modeling choices,
not a model-free repair for weak or invalid instruments.
\citep{exec25_pyrdiv_causalpy}

A bank request must identify the coefficient associated with the
instrumented treatment, the structural interpretation, instrument and
outcome populations, row identities, priors and sampling diagnostics.
Formula construction and treatment-name checks cannot prove exclusion,
monotonicity, truthful encouragement or a stable target population.
Data-centered priors should be distinguished from externally fixed priors
when interpreting sensitivity and uncertainty. Variable selection is
another model operation requiring its own contract, rather than evidence
that automatically discovered instruments are valid.
\citep{exec25_pyrdiv_causalpy,exec25_pyrdiv_choice}

The proposed use is a version-pinned model-based reference and conditional
Python worker. Public tests should check formula alignment, transformed
treatments, repeated fitting, finite inputs, weak first-stage scenarios,
prior sensitivity and posterior diagnostics before native admission.
None of those tests can manufacture the banking facts that make an
instrument credible.~\citep{exec25_pyrdiv_choice}

\subsection{A finite cutoff prediction is not automatically the limiting jump}
The selected CausalPy RD implementation fits a sharp-design regression
after optional bandwidth and donut-hole filtering. It predicts at the
threshold minus and plus a user-supplied epsilon, with treatment flags zero
and one, and subtracts the two conditional expectations. Its input checks
require a dummy-coded treatment column and a formula treatment term;
the inspected checks do not establish that the observed assignments match
the declared cutoff direction. The adapter must validate that relationship
and preserve the actual prediction rows.~\citep{exec25_pyrdiv_causalpy}

Let a local mean below the cutoff have slope 2 and the mean above it have
slope 6, with a limiting jump of 10. Comparing symmetric points at distance
$\epsilon$ gives $10+8\epsilon$. The contrast is 14 at $\epsilon=1/2$ and
10.008 at $\epsilon=0.001$. These authored values show a target difference,
not numerical error in a package. The units and transformation of a credit
score affect the meaning of epsilon. A donut design adds extrapolation
across an excluded region; its target and assumptions must remain explicit.
\citep{exec25_pyrdiv_choice}

The small rdd 0.0.3 package provides a different contract. Its helper can
truncate data around a cutoff and another function creates a treatment
indicator and a statsmodels weighted-least-squares model. The default
formula includes the treatment indicator and the running variable with a
common slope; the function returns the model before fitting. It does not
automatically perform the separate truncation step or turn the default
formula into a slope-interacted local polynomial. Users can supply a
different equation and weights, but those choices and the fitted covariance
must be recorded.~\citep{exec25_pyrdiv_rdd}

For the proposal, rdd is a compact instructional/reference route. Direct
adoption for a bank RD workflow is deferred until a wrapper fixes window,
kernel, formula, fitted state, covariance and diagnostics. Its MIT notice
does not qualify the complete statsmodels environment, and a coefficient
labeled as an effect still needs the design's causal assumptions.
\citep{exec25_pyrdiv_rdd,exec25_pyrdiv_choice}

\subsection{rdrobust keeps bias correction and its uncertainty in separate rows}
The pinned rdrobust Python source declares version 2.1.1 and GPL-3 terms.
Its interface distinguishes sharp and fuzzy designs, polynomial and bias
orders, derivative targets, kernels, bandwidths, covariates, weights and
variance options. The inspected guards require the bias order to exceed
the estimation order, reject negative or all-zero weights, validate a
finite interior cutoff and require enough distinct running-variable values
for the polynomial. The treatment-side partition places values at the
cutoff on the right. Repeated-score handling and numerical refusal are
useful operational checks, not proofs of the continuity design.
\citep{exec25_pyrdiv_rdrobust}

The returned rows pair conventional and bias-corrected estimates with
different standard errors. Specifically, the selected construction uses
the conventional estimate with the conventional standard error, the
bias-corrected estimate with that conventional standard error, and the
bias-corrected estimate with its robust standard error. A consumer must
preserve the row identity; extracting the corrected coefficient from one
row and a standard error from another changes the reported test.
\citep{exec25_pyrdiv_rdrobust}

Suppose an authored example has conventional estimate 10, corrected
estimate 8, conventional standard error 2 and robust standard error 3.
The three reported ratios are 5, 4 and $8/3$. With a stipulated normal
multiplier 1.96, the corresponding interval half-widths are 3.92, 3.92
and 5.88. No coverage claim follows from this arithmetic. It provides a
deterministic serialization oracle that prevents a bridge from silently
mixing rows.~\citep{exec25_pyrdiv_choice}

For fuzzy designs the inspected formulas divide an outcome discontinuity
by a treatment discontinuity and propagate the ratio through bias and
variance calculations. A guard rejects a fully degenerate treatment
variable with no variation and no jump. That guard does not establish a
strong first stage. An outcome jump of two divided by treatment jumps
0.2 and 0.02 yields 10 and 100 respectively. A narrow conventional interval
cannot be assumed reliable merely because division returned a finite
number. Weak-first-stage behavior and the intended local population need
separate scientific evaluation.~\citep{exec25_pyrdiv_rdrobust,
exec25_pyrdiv_choice}

The proposal favors a supported rdrobust Python worker as a reference
candidate for declared local-polynomial tasks, conditional on exact artifact
identity, complete bandwidth/covariance delegates, license and dependency
review, and matched native public fixtures. It does not port the method
solely to avoid process isolation, nor claim general serial-dependence or
weak-instrument coverage from an option label.
\citep{exec25_pyrdiv_choice}

\subsection{ivmodels makes confidence-set shape part of the API}
The selected ivmodels 0.10.0 source offers coefficient-null Anderson--Rubin
tests and inversions. Its inspected path residualizes controls, projects
null residuals onto instruments, constructs a statistic and selects an
explicit chi-square, F or supported subvector critical-value route.
The inversion forms a quadratic inequality and projects it onto requested
coefficient coordinates. This is distinct from the earlier relevance and
overidentification diagnostics. No general covariance-robust guarantee is
inferred from the method name or source comments.
\citep{exec25_pyrdiv_ivmodels}

The confidence-set class represents unions of intervals, including empty,
bounded and unbounded sets. That representation is substantively important.
The authored inequality $1-\beta^2\leq0$ accepts
$(-\infty,-1]\cup[1,\infty)$ and excludes zero. Replacing it by the
interval spanning its extremes would add excluded parameter values.
A banking result schema must represent disconnected, empty and unbounded
sets, together with the null, nuisance treatment, critical-value option,
degrees of freedom and numerical tolerance.~\citep{exec25_pyrdiv_ivmodels,
exec25_pyrdiv_choice}

The numerical boundaries still require work. The inspected projection
delegate uses SciPy least squares while the quadratic constructor obtains
its center with a direct linear solve. Singular or nearly singular
quadratic boundaries, numerical rank, root search and complete caller
validation need explicit fixtures. An algebraic acceptance region can
exist when a particular center parameterization is singular. Native
qualification must distinguish a scientific empty set from a numerical
failure. The proposed role is a weak-IV inference reference and conditional
worker integration preserving set geometry, not automatic admission of
all instrument laws or dependence regimes.~\citep{exec25_pyrdiv_ivmodels,
exec25_pyrdiv_choice}

\subsection{Public qualification before costly banking experiments}
The public package should test exact cutoff inclusion, assignment direction,
epsilon and score scaling, donut filtering, formula slopes, duplicate rows,
finite weights, rank tolerance, solver status, weak first stages, robust-row
serialization and disconnected confidence sets. It should distinguish
simulation under a declared valid law from deliberate violations of
exclusion, continuity or dependence assumptions. Costs include refusals,
failed requests and human review, not only returned estimates.
\citep{exec25_pyrdiv_choice}

A hypothetical 18,000 USD total over 600 attempts, 450 conformant results
and 300 useful results gives 30, 40 and 60 USD respectively. These are
planning denominators, not observed prices or banking value. Repository
search tails, nonlinear IV alternatives, complete dependence and
bandwidth delegates, artifact rights and native evidence remain open.
The comparison of Python instruments and discontinuities remains incomplete.
The packages have not been validated by execution in this comparison.
\citep{exec25_pyrdiv_choice}
